\documentclass[../main.tex]{subfiles}
\begin{document}
%==========================================TIÊU ĐỀ TẬP========================================
\begin{table}[h]
    \begin{adjustwidth}{-1cm}{}
        \begin{tabular}{>{\centering\arraybackslash}p{6cm}|>{\raggedright\arraybackslash}p{12.5cm}}
            \multirow{5}{6cm}
            % Cột bên trái: ÉPISODE và số tập
            {\\[-40pt]
                \flaregothic\fontsize{20pt}{20pt}\selectfont \centering
                \color{eptype}ÉPISODE\color{black}\\
                \burbank\fontsize{72pt}{72pt}\selectfont
                \color{epnum}62\color{black}\\[6pt]
                \cafeteria\fontsize{20pt}{20pt}\selectfont\color{epnum}(5-6)\color{black}
                \\[18pt]
                \flaregothic\fontsize{14pt}{14pt}\selectfont
                \color{epattr}BIENVENUE, \uline{\color{PPTenshi} Kanata} !
                \color{black}
                \\}
             &
            % Dòng thứ nhất: tên tập tiếng Nhật
            {\fotskip\fontsize{18pt}{18pt}\selectfont \ruby{関}{かん}\ruby{係}{けい}が\ruby{ギクシャク}{ぎくしゃく}する\ruby{爆}{ばく}\ruby{彈}{だん}\ruby{解}{かい}\ruby{除}{じょ}
            } \\[6pt]
            % Dòng thứ hai: tên tập tiếng Anh
             & {\cafeteria\fontsize{18pt}{18pt}\selectfont Bomb Disposal Squad}                                                                     \\[4pt]
            % Dòng thứ ba: tên tập tiếng Việt
             & {\vnmsans\fontsize{18pt}{18pt}\selectfont Gỡ bom}                                                                                    \\[8pt]
            % Dòng thứ tư: Nhân vật xuất hiện
             & {\caviar{\fontsize{14pt}{14pt}\selectfont Nhân vật: \emp{kuon1} \emp{ayame} \emp{choco} \emp{coco} \emp{kanata}}}
            \\[8pt]
            % Dòng cuối cùng: Ngày phát hành
             & {\continuum\fontsize{16pt}{16pt}\selectfont Ngày phát hành: 12/07/2020}                                                              \\[18pt]
        \end{tabular}
    \end{adjustwidth}
\end{table}
%=========================================TRUYỆN CHÍNH========================================
\color{black}

% Hộp tóm tắt
\txtbx[2]{{\color{vol} \uline{Tóm tắt:}} Văn phòng \hll\ xuất hiện một quả bom - thủ phạm là Ayame vì sợ gián. Choco lập tức thành lập ``Đội gỡ bom truyền kỳ'' với Coco, Kanata và Kuon. Trong lúc căng thẳng, Coco cắt đại dây đỏ (dù Kuon khuyên cắt xanh), ném bom vào tủ lạnh, và cuối cùng vẫn nổ tung. Vụ nổ thiêu rụi văn phòng, bốn cô gái rơi từ trên cao xuống đất đen thui. Kuon gọi cho người bạn Koheki để than thở, và phát hiện mình đã bị lừa từ đầu.}

\begin{center}
    \uline{(Người viết tập này có sử dụng bản dịch từ \href{https://www.youtube.com/@TomangclipHololivevedich}{Tớ mang clip \hll\ về dịch - TMCHVD}.)}
\end{center}

\pov{Hikawa Kuon}

Tuần thứ hai, chúng ta có ai sẽ đến văn phòng đây?

Một thiên sứ đến từ Học viện Thiên đàng, đang phấn đấu để trở thành một thiên sứ chính thức - hay tôi còn nói vui là "thiên sứ xịn". Tôi bảo thế bởi vì trên tay em ấy vẫn còn đeo băng "Thiên sứ học viện", chứng tỏ rằng vẫn chưa tốt nghiệp hoàn toàn. Có lẽ em ấy đang học năm cuối, kiểu như sinh viên thực tập vậy.

Cô nàng này có một tinh thần thép - dù vẻ ngoài trông như một thiên sứ hiền lành, nhưng thực chất lại là một chiến binh. Nghe nói lực tay của em lên tới năm mươi cân, đủ để bóp nát một quả táo chỉ bằng một tay. Bền bỉ, chăm chỉ và đầy quyết tâm, em ấy luôn cố gắng hết sức để vươn lên trong thế giới này - có lẽ vì cánh cổng thiên đàng không dành cho những kẻ lười biếng.

Tuy nhiên, tính cách của em lại khá thú vị: là một người vui vẻ, hòa đồng khi ở bên những người thân thiết, ví dụ như Coco; nhưng lại có tính cách lạnh lùng (thực ra là nhút nhát) và ít nói khi đối mặt với người lạ - tôi, chẳng hạn. Lần đầu gặp, em ấy gần như không nói câu nào, chỉ gật đầu và cười nhẹ. Không giỏi giao tiếp với đám đông, nhưng bù lại, em ấy có một giọng hát ấn tượng - đến mức tôi từng ngỡ rằng thiên đàng đang hát qua loa.

Em ấy cũng là một chuyên gia PowerPoint, và tôi không đùa. Có lần tôi thấy em ấy làm một bản trình chiếu về cách làm bánh, và nó đẹp đến mức tôi nghĩ em ấy nên xin việc ở công ty thiết kế đồ họa. Ngoài ra, em ấy còn có tài nhại giọng người khác - một kỹ năng mà tôi nghĩ sẽ rất hữu ích khi muốn trêu chọc đồng nghiệp.

Và dù là một thiên sứ, mục tiêu lớn nhất của em lại không hẳn là làm một "sứ giả thiên đàng". Em ấy muốn một ngày nào đó tổ chức một buổi hòa nhạc tại Saitama Super Arena. Một thiên sứ với một giấc mơ đầy tham vọng.

\dots Hoặc có lẽ, một thiên sứ sa ngã thích nổi loạn. Tôi không biết thiên đàng có cho phép thiên sứ hát rock ở Saitama Arena không, nhưng tôi nghĩ là có - miễn là họ không phá hỏng thiết bị.

\chrint

Đó là Amane Kanata (\ruby{天}{あま}\ruby{音}{ね}かなた, \ruby{A-ma}{Thiên}-\ruby{nê}{Âm}\ Ca-na-ta), bạn thân của Coco, và cũng là thành viên thứ hai của \textit{holoForce}.

Một cô gái 18 tuổi (thiên thần, nhưng em ấy nói rằng tuổi này bằng tuổi tính theo con người, và tôi không hỏi thêm vì sợ bị sét đánh), và như đã nói ở trên, là một thiên thần học việc tại Học viện Thiên đường. Em ấy vẫn chưa có cánh lớn, chỉ có một đôi cánh nhỏ xíu trên lưng, trông như đồ trang trí hơn là thứ để bay.

Kana-chan cực kỳ ngọt ngào, lịch sự, thân thiện, hơi ngốc nghếch, trong sáng và khiêm tốn. Em ấy có phần trầm tính và kín đáo hơn khi so sánh với những người cùng Thế hệ thứ tư, và khiếu hài hước của em ấy thường có xu hướng ít thô tục hoặc khiếm nhã hơn, trái ngược hoàn toàn với bạn thân Coco - người mà câu chuyện cười thường bắt đầu bằng một tiếng chửi thề bằng tiếng Anh.

\begin{wrapfigure}[21]{l}{6cm}
    \includegraphics[width=8cm]{../../../../Image/Book5kanata_model.png}
\end{wrapfigure}

Hòa nhã và vui vẻ là thế, Kana-chan lại có một tính cách hoàn toàn ngược so với một thiên thần truyền thống: một "ác quỷ" tiềm ẩn. Bằng chứng là cách em ấy trêu chọc khán giả hoặc chơi khăm những thành viên \hll\ khác. Có lần em ấy đã giả vờ khóc để lừa một đồng nghiệp, và khi người đó chạy đến an ủi, em ấy bật cười và nói "Đùa thôi!" - một kiểu hài hước khiến tôi nhớ đến một số người khác trong công ty.

Kana-chan được người xem và đồng nghiệp nhận xét là thực sự có phần ngây thơ, thậm chí "trong sáng" đến mức đáng lo - nhiều người hâm mộ thường đùa rằng cô ấy có thể bị "làm hư hỏng" bởi những thành viên khác, đặc biệt là Coco. Tôi nghĩ chính sự tương phản này làm cho em ấy trở nên thú vị.

Như đã trình bày ở trên, em ấy là một người có giọng hát được hai bố con tôi - những người cảm thụ âm nhạc - đánh giá là trời phú. Tôi nhớ lần đầu nghe em ấy hát, tôi đã nghĩ: "Thế này mà gọi là thiên sứ học việc à? Chắc giáo viên dạy nhạc ở thiên đàng phải là dân chuyên nghiệp." Độ ngân nốt cao của em ấy có thể lên đến hai mươi giây, một mức mà rất ít người làm được, kể cả các ca sĩ chuyên nghiệp.

Kana-chan cũng là một người có nghị lực rất lớn. Em ấy từng phải đối mặt với những khó khăn, như bị điếc một bên tai, nhưng vẫn kiên trì theo đuổi ước mơ ca hát. Tôi nghĩ rằng, chính những thử thách đó đã tôi luyện nên một tinh thần thép bên trong vẻ ngoài mong manh của một thiên sứ.

Về vẻ bên ngoài, Kana-chan có đôi mắt màu tím-hồng, mái tóc ngắn màu xám với phần tóc bên trong màu xanh lam và một vài vệt màu hồng. Em ấy để kiểu tóc bob bất đối xứng (với một lọn tóc đơn và mái dài bất đối xứng che mắt phải) - kiểu tóc khiến người ta muốn vuốt, nhưng tôi không dám vì sợ bị đập năm mươi cân. Em cũng có đôi cánh thiên thần nhỏ rời mang hai màu từ trắng sang xanh lam, và một vầng hào quang hình ngôi sao (đôi khi được so sánh với shuriken - phi tiêu) lơ lửng lệch tâm trên đầu. Tôi đã thấy em ấy dùng hào quang đó để cắt băng keo một lần, và từ đó tôi không còn coi nó là vật trang trí nữa.

Nàng Thiên sứ này mặc một chiếc áo khoác ngắn màu xám với một chiếc áo có tay rộng và dài qua cổ tay với họa tiết ren và cổ áo thủy thủ màu trắng. Bên ngoài là áo sơ mi trắng và váy cổ lọ màu đen (mặc dù thường được mô tả là chỉ là một chiếc áo cổ lọ riêng biệt và váy ngắn xếp nếp màu đen), đeo một chiếc nơ màu xanh trước ngực, một chiếc băng tay màu vàng đề ba chữ "生徒会 - Sinh Đồ Hội, tức là hội đồng học sinh". Em ấy đi tất ca-rô màu xanh-trắng dài đến gần đầu gối, kèm theo dây nịt quanh cổ chân và bốt đen. Trông em ấy như một học sinh gương mẫu, nhưng tôi đã học được rằng đừng bao giờ đánh giá một thiên sứ qua vẻ ngoài.

\chrint

Để xem nàng Thiên sứ này sẽ bày trò gì cho văn phòng \hll\ chúng ta nào.

\il

Tuần đầu tiên với Coco đã là một cơn ác mộng đầy màu sắc. Tôi đã chứng kiến một cô Rồng biến văn phòng thành sân khấu đột kích, phá cửa kính, cưỡi khủng long, và thề thốt bằng tiếng Anh mỗi năm phút. Tôi tưởng rằng tuần thứ hai sẽ bình yên hơn. Tôi đã sai. Và tôi sai một cách ngoạn mục.

Ngay khi tôi vừa bước vào văn phòng, chưa kịp ngồi xuống nghỉ ngơi sau một buổi sáng mệt mỏi, thì nàng Y tá Yuzuki Choco bước vào với một tin sốc làm tôi quên luôn cái mệt: ``Mọi người biết không? Có bom trong văn phòng này đấy.''

Tôi quay ngoắt lại, tay vẫn còn cầm tách cà phê, mắt trợn lên như hai viên bi: ``Cái gì cơ? Bom? Bom thật hay bom giả? Hay là trò đùa của ai đó?''

Coco, người cũng vừa mới tới cùng tôi từ hành lang, cũng sửng sốt không kém. Mắt cô mở to, cái miệng vốn hay chửi thề của cô lần này chỉ kịp thốt lên: ``Chúng ta có bom à? Ở văn phòng này? Ai mà có gan đặt bom ở đây?''

Choco gật đầu, khoanh tay lại, giương ra bộ mặt lo lắng đầy kịch tính. Cô ấy có vẻ nghiêm túc, nhưng cũng có vẻ như đang tận hưởng cảm giác làm nhân vật chính trong một bộ phim hành động: ``Phải. Có vẻ như ai đó hoảng quá sau khi thấy mấy con gián lảng vảng trong này, thế là đặt vài quả bom để `xử lý triệt để'. Họ muốn diệt gián, nhưng diệt luôn cả văn phòng.''

Tôi xoa trán, hơi trố mắt, giọng đầy bất lực: ``\vie{Ùi, vãi.} Liều đến mức ấy cơ à? Diệt gián bằng bom thì sau khi nổ, không còn văn phòng, không còn gián, cũng không còn chúng ta. Một mũi tên trúng ba đích.''

Choco và tôi lo lắng là vậy, nhưng với Coco - người không thiếu những trò nghịch ngợm - có vẻ hào hứng hơn hẳn. Cô nhảy lên một cái, hai tay giơ cao như thể vừa nhận được tin về một trò chơi mới: ``Chơi khô máu với đám gián luôn nhỉ. Ai mà nghĩ ra được kế này thì đúng là liều mạng thật.''

Tôi thì chẳng thấy vui chút nào, tôi nhìn quanh văn phòng, mắt dõi theo từng góc tủ, từng kệ sách: ``Thế này thì gián chưa chết mà chúng ta đã mất cả đống tiền rồi đấy\dots\ Sửa lại văn phòng còn đắt hơn mua thuốc diệt gián gấp trăm lần.''

Ngay lúc đó, một giọng nói hào hứng vang lên phía sau, như thể một con chim nhỏ vừa đáp xuống cành: ``Ồ, có chuyện gì vậy? Có bom à? Tôi nghe nói có bom!''

Tôi liếc xuống, thấy Kanata - nàng Thiên sứ nhỏ - đang nhảy lên trên cao rồi đáp xuống, vầng hào quang trên đầu xoay tít như một chiếc đĩa bay mini.

\gif{5-6a}{1}{42}

Coco nhướng mày, lẩm bẩm với người đồng đội của mình, giọng vừa nghi ngờ vừa thích thú: ``Đứa dở hơi nào lại mang chất nổ vào đây vậy? Mà cũng hay, ít nhất cũng làm không khí sôi động hơn.''

Tôi phỏng đoán ngay, trong đầu lần lượt điểm danh các ứng cử viên: ``Khả năng cao là Haachama hoặc Peko-chan. Hai đứa đó sợ gián nhất mà. Mà có khi là cả hai, chúng nó thường làm mấy trò điên rồ cùng nhau.''

Kanata đưa tay chỉ về phía kệ tủ gần đó, làm cả ba chúng tôi cùng quay đầu nhìn theo: ``Có phải là cái mà mọi người đang nói đến phải không? Trông nó giống bom thật đấy.''

Choco gật đầu, giọng đầy trách nhiệm của một chỉ huy: ``Đúng rồi đấy. Đó là mục tiêu của chúng ta.''

Tôi thì chỉ có một thắc mắc khác, một thắc mắc đơn giản nhưng đầy nghi vấn: ``Ai lại đặt bom ở chỗ dễ thấy vậy nhỉ? Nếu là tôi, tôi sẽ giấu nó dưới gầm bàn hoặc trong tủ hồ sơ. Đặt trên kệ như trưng bày thì trông lộ thiên quá.''

Nhưng thứ thu hút sự chú ý của tôi hơn cả là hai con búp bê Matryoshka bằng sứ đặt trên kệ. Chúng đang nói chuyện với nhau, bằng một cách nào đó, và tôi không nghĩ đó là ảo giác của mình.

\img{5-6b}

``Chà, chỗ này ngồi thoải mái thật nhỉ,''  con búp bê mang dáng dấp Kanata nói, giọng y hệt, đến mức tôi suýt quay sang hỏi Kanata thật.

Con búp bê còn lại, có ngoại hình giống Ayame, nhìn về phía chúng tôi, liền giật bắn mình nhắc đồng đội, giọng đầy hoảng loạn: ``A(n)dy kìa\footnote{Liên hệ từ \textit{Toy Story} (Câu chuyện đồ chơi).}! Mọi người, mau giả làm đồ chơi nào! Họ đang nhìn kìa!''

Ngay lúc đó, con búp bê màu đỏ ngả người về phía sau, vô tình chạm vào một sợi dây. Một tiếng ``tích'' vang lên, rồi một màn hình số nhỏ bật sáng trên quả bom gắn trên kệ.

Bốn phút đếm ngược bắt đầu.

Cả bốn chúng tôi sững sờ. Không ai nói gì trong một giây đầu tiên. Rồi mọi người cùng nhìn nhau, mắt mở to, miệng há hốc. Tôi chưa kịp phản ứng gì thì chỉ kịp thốt lên một câu, giọng đầy bất lực: ``\eng{\dots Chết tiệt.}''

\dots

\dots

Ngay sau đó, Kanata hoảng hốt xoay tròn cả hào quang trên đầu như cái chong chóng, vầng sáng xoay tít tạo thành một vòng tròn mờ ảo. Còn Coco thì gần như reo lên vì phấn khích - cũng phải thôi, dường như em ấy có cớ để bày thêm trò mới rồi: ``Chúng ta cần phải mau gỡ quả bom này! Đây là nhiệm vụ đầu tiên của holoForce! Tuyệt vời!''

Choco theo đó, nhanh chóng đưa ra hiệu lệnh, giọng như một vị tướng trên chiến trường: ``Đội gỡ bom, tập hợp! Chúng ta có bốn phút để cứu văn phòng!''

Chưa đầy một phút sau, cả ba cô gái lẫn tôi (một cách bất đắc dĩ) đã đội sẵn\dots\ mũ bảo hộ lao động. Trông chúng tôi như một đội phá dỡ nhà cửa chứ không phải một đội gỡ bom.

Choco với phong thái nghiêm túc, chỉ tay về quả bom trong kệ đó: ``Nhiệm vụ: Gỡ bom trong vòng bốn phút! Mọi người, đã rõ chưa!?''

Hai thành viên của thế hệ thứ tư giơ tay chào kiểu nhà binh, giọng vang dội: ``Yes, sir! (Đã rõ!)''

\img{5-6c}

Tôi nhìn hai người họ, chớp mắt\dots\ rồi cũng giơ tay chào theo mà chẳng hiểu sao mình lại bị cuốn vào vụ này. Cánh tay tôi giơ lên, miệng thốt ra một câu gì đó mà tôi không nhớ, nhưng trong đầu tôi chỉ có một suy nghĩ duy nhất: ``\textit{Thật sự mình đang đội mũ bảo hộ, chuẩn bị gỡ bom, trong một văn phòng, vì mấy con gián ư?}''

\ct

Trưởng đội - Choco - đang nấp phía sau một đống thùng các-tông, giữ khoảng cách an toàn với hiện trường. Cô cầm bộ đàm, bình tĩnh thông báo, giọng như một chỉ huy thực thụ trong phim chiến tranh: ``Chúng ta còn ba phút đến khi quả bom phát nổ! Mọi người nhanh tay lên đi ha.''

Ở tiền tuyến, ba kẻ xấu số là Coco, Kanata và tôi đang cố tìm cách ngăn cản cỗ máy hủy diệt này. Vấn đề là\dots\ không ai trong ba chúng tôi có kinh nghiệm xử lý bom cả. Tôi đã từng sửa máy tính, từng xử lý đống giấy tờ, từng dàn xếp mấy vụ cãi vã của các cô gái, nhưng gỡ bom thì chưa bao giờ.

Trong khi tôi và nàng Rồng còn đang nhăn nhó suy nghĩ, nàng Thiên sứ tóc trắng - người được cho là đáng sợ nhất trong tình huống này - lại tỏ ra khá thoải mái. Không biết làm gì cho bớt nhàm chán, cô đưa ra một ý tưởng: ``Ta có thể làm ra\dots''

Nhưng chưa kịp nói hết câu, Coco liền quay ngoắt sang, ánh mắt sắc bén như muốn kết liễu ai đó ngay lập tức: ``Nếu cô dám nói `ramen' thì Kanatan sẽ là đứa tiếp theo đấy!''

Em ấy vừa nói vừa túm lấy con búp bê sứ màu xanh bên cạnh và ném thẳng xuống đất, khiến nó vỡ vụn thành hàng tá mảnh nhỏ. Nhưng kỳ lạ thay, nó được phục hồi quay trở về hình dạng ban đầu, như một phép màu. Có phải bởi vì Kanatan là thiên sứ chăng?

Kanata đứng hình. Cô im lặng mất vài giây, miệng há hốc, rồi lại định lên tiếng: ``Cái này cứ như là trong U--''

Lần này, Coco không đợi cô nói hết câu: ``Và nếu định nói cái này giống như `Ultraman' thì cô sẽ là đứa tiếp theo đấy!''

Nói rồi, cô lại túm lấy con búp bê sứ còn lại - con có hình dáng Ayame - và bóp vụn nó trong tay như đang nghiền nát vận mệnh của một người nào đó. Khác với con búp bê trước, con này đã tan tành, không quay trở về hình dáng cũ được nữa. Những mảnh vụn rơi xuống sàn, tạo thành một vũng sứ vỡ. Có lẽ bởi vì Ojou là quỷ Oni - kiểu mục tiêu đáng bị tiêu diệt - rồi.

\begin{figure}[H]
    \centering
    \begin{subfigure}[t]{0.45\textwidth}
        \centering
        \animategraphics[width=8cm]
        {30}{../../../../Image/Book5/5-6e1/5-6e1.}{1}{84}
        \caption{Coco ném búp bê Matryoshka hình Kanata, nó lập tức hồi về hình dáng cũ}
    \end{subfigure}
    \quad
    \begin{subfigure}[t]{0.45\textwidth}
        \centering
        \animategraphics[width=8cm]
        {30}{../../../../Image/Book5/5-6e2/5-6e2.}{1}{54}
        \caption{Coco làm tương tự với búp bê hình Ayame, nhưng nó vỡ vụn.}
    \end{subfigure}
\end{figure}

Tôi rùng mình, vội ho nhẹ một tiếng để phá tan bầu không khí căng thẳng. Tôi nhìn Coco, rồi nhìn Kanata, rồi nhìn đống búp bê vỡ: ``A-Anh nghĩ là chúng ta nên tập trung vào gỡ bom đi em ạ\dots\ Có gì cứ nói sau\dots\ Nếu tiếp tục thế này, chúng ta sẽ hết búp bê để phá trước khi tìm ra dây.''

Tôi liếc sang nàng Thiên sứ, mắt em ấy trơ ra, như thể vừa bị dọa đến mức không còn cảm xúc. Cô nhìn nàng Rồng, rồi nhìn tôi, rồi nhìn lại Coco, miệng mấp máy nhưng không phát ra âm thanh.

Bất ngờ thay, Coco không làm gì tôi cả. Có lẽ bởi vì tôi là yếu nhân, hoặc là em ấy biết rõ tôi là Tổng quản lý nên không làm càn. Có thể em ấy nghĩ rằng bóp vỡ tôi sẽ không có ai chịu trách nhiệm sửa chữa.

\dots Thôi kệ. Tập trung vào gỡ bom nào.

\ct

Sau khi bình tĩnh lại, ba người quay lại với công việc. Nhưng bộ đàm lại réo lên với giọng của trưởng đội, vang lên từ phía sau đống thùng: ``Tình hình ở đó thế nào rồi? Hết!''

Kanata lập tức bấm bộ đàm, nhưng thay vì báo cáo tình trạng của quả bom, em ấy lại nói bằng giọng sợ sệt, như một đứa trẻ vừa thấy ma: ``Đội gỡ bom đây, cô ta lạnh nhạt quá, hết!''

Tôi chưng hửng, giật lấy bộ đàm từ tay cô, giọng đầy bất lực: ``Báo cáo, bọn anh đang tháo dỡ vỏ bom để tìm dây. Hết!''

Bên kia bộ đàm, giọng Choco đáp lại, đầy phấn khích của một chỉ huy được cập nhật thông tin: ``Hiểu rồi. Cố lên, các em! Hết.''

Ngay sau đó, Kanata bỗng nhiên hào hứng reo lên, tạo dáng như một siêu anh hùng trong phim hành động, một tay giơ lên trời, một tay đặt lên ngực: ``\textbf{Đội gỡ bom truyền kỳ, xuất kích!}''

\img{5-6e}

Cô giữ nguyên tư thế trong vài giây, nhưng không có câu trả lời từ phía bên kia đầu dây. Không khí im lặng đến đáng sợ.

``Cô ấy tắt máy rồi!''  cô kêu lên, tụt hứng như một diễn viên vừa bị cắt mic giữa sân khấu.

Tôi liếc nhìn em ấy tạo dáng, ngơ ngác thắc mắc với vẻ mặt đầy hoài nghi: ``\dots Em đang làm cái trò gì vậy? Đội gỡ bom truyền kỳ à? Anh chưa từng nghe nói đến tên đó bao giờ.''

Kanata nghe tôi nói vậy liền đỏ mặt, suýt vung tay vào người tôi, giọng lắp bắp: ``K-Không có gì! Chỉ là\dots\ em nghĩ rằng cần có một cái tên để tạo khí thế thôi.''

\ct

Dù hơi sốc vì trò cười mà em ấy bày ra, chúng tôi vẫn tiếp tục với công việc. Và rồi, Coco phát hiện ra một thứ quan trọng - thứ có thể là chìa khóa để cứu cả văn phòng. Em ấy ngồi xuống, cúi người nhìn kỹ vào bên trong quả bom: ``Ê, tôi tìm được mấy sợi dây điện rồi này!''

Trước mắt chúng tôi là hai sợi dây, một xanh, một đỏ, được bố trí gọn gàng nhưng đầy đe dọa. Mỗi sợi đều có một ký hiệu nhỏ, nhưng tôi không đọc được vì ánh sáng mờ.

Kanata nghe vậy, mắt sáng rực lên như thể vừa trúng số độc đắc. Cô ấy vỗ tay, nhảy dựng lên, hào quang trên đầu xoay tít vì phấn khích: ``A, Conan có dạy tôi về cái này rồi nè! Trong tập mà cậu ấy giải cứu cả một đoàn tàu!''

\img{5-6f}

Tôi vội quay sang, tràn đầy hy vọng. Có lẽ vị cứu tinh đã đến. Tôi mở to mắt, giọng như cầu cứu: ``Ủa, thật à? Vậy em có biết ta nên làm gì không? Cắt sợi nào? Xanh hay đỏ?''

Kanata gật đầu, tỏ vẻ chắc nịch, như một chuyên gia xử lý bom đã có kinh nghiệm: ``Ta cứ cắt dây thôi!''

Câu trả lời của em ấy khiến tôi ngã ngửa suýt đập đầu vào quả bom: ``Cái đấy thì ai chả biết!? Cắt dây, nhưng cắt dây nào mới là vấn đề! Mà con mắt siêu hạng, kinh nghiệm gỡ bom của em đâu?''

Coco gãi cằm, lẩm bẩm, mắt nhìn hai sợi dây như thể đang soi một ván cờ: ``\dots ta phải cắt sợi nào đây? Xanh hay đỏ? Trong phim họ hay cắt đỏ, nhưng cũng có tập cắt xanh\dots''

Tôi nuốt khan, nhìn hai sợi dây đó với đồng hồ đang đếm ngược dần về hai phút. Mồ hôi bắt đầu lấm tấm trên trán: ``Năm mươi năm mươi đấy. Theo kinh nghiệm của anh thì\dots''

Tôi dừng lại, vì thực ra tôi chẳng có kinh nghiệm gì. Nhưng tôi không thể để họ biết điều đó.

``\dots ờm, anh nghĩ là nên cắt sợi màu đỏ. Vì\dots\ màu xanh là màu của thiên nhiên, của hòa bình. Còn màu đỏ là màu của báo động, của nguy hiểm.''

Kanata nhìn tôi với ánh mắt đầy ngưỡng mộ: ``Anh đúng là có kinh nghiệm thật!''

Coco chỉ nhún vai, giơ kéo lên: ``Thôi, cắt xanh đi. Nếu sai thì chúng ta sẽ bay lên thiên đàng cùng nhau.''

Tôi nhìn em ấy, mồ hôi túa ra nhiều hơn: ``Em có thể đừng nói câu đó được không? Nó không giúp ích gì cho tinh thần đâu.''

\ct

\pov{-}

Trong lúc tiền tuyến đang căng như dây đàn, ở hậu phương - đằng sau những chiếc hộp bìa các-tông - vẫn hoàn toàn yên tĩnh. Quá yên tĩnh, đến mức đáng ngờ. Không một tiếng động, không một tiếng bước chân, chỉ có tiếng ve kêu râm ran từ bên ngoài.

Nàng Y tá vừa ngồi chờ trong sự chán nản, vừa lẩm bẩm một mình, ngón tay gõ nhịp lên bộ đàm: ``Không biết ba người họ có ổn không nhỉ\dots\ Đã hơn một phút rồi mà chưa có tin tức gì. Có khi nào họ bế tắc điều gì đó sao?''

Cô đang ngồi ôm bộ đàm, chờ tin tức từ nhóm gỡ bom thì bất chợt nhìn thấy một bóng dáng quen thuộc đang run run dưới gầm bàn. Một mái tóc tráng cột búi, một đôi sừng nhỏ, đang ngồi xổm xuống, hai tay ở đằng sau gáy, lưng run lên bần bật.

\img{5-6g}

Thấy làm lạ, Choco cúi xuống, giọng đầy ngạc nhiên: ``Em trốn trong đó làm gì vậy? Có chuyện gì mà phải chui xuống gầm bàn thế?''

Ayame ngẩng lên, mặt tái mét như vừa nhìn thấy một thế lực hắc ám nào đó. Hoặc chính xác hơn, cô đang ngỡ ngàng trước một sự thật khủng khiếp. Giọng cô run run, đầy thương tiếc: ``Ai đó đã phá mấy món đồ chơi gia truyền bị nguyền rủa của gia tộc em rồi\dots\ Chúng đã tồn tại qua bao thế hệ, là bảo vật của dòng họ Nakiri\dots''

Nói rồi, nàng hít một hơi thật sâu, giọng run rẩy tiếp tục, như một người vừa nhận tin dữ: ``Giờ chỉ còn lại mấy quả bom thôi! Tất cả những gì em có đều đã tan tành!''

Choco đơ người mất vài giây, chưa kịp hiểu ra vấn đề. Cô nghiêng đầu, giọng đầy thắc mắc: ``Em đang nói cái gì vậy? Đồ chơi gì? Bom gì? Em có thể nói rõ hơn được không?''

Chưa kịp nhận câu trả lời, Ayame bỗng hét lên đau đớn, giọng vang vọng khắp căn phòng, xé toạc bầu không khí im lặng: ``\textbf{A(n)dy!!}''

Tiếng hét của nàng Quỷ vọng ra tận hành lang, khiến cả những người ở tiền tuyến cũng phải giật mình. Choco hoảng loạn, vội vã chui xuống gầm bàn, cố gắng dỗ dành nàng Quỷ như đang trấn an một đứa trẻ vừa làm vỡ đồ chơi yêu thích: ``Nào nào! Bình tĩnh lại đi chứ! Chỉ là đồ chơi thôi mà, không có gì phải sợ cả! Em có thể mua lại được mà!''

Nhưng Ayame vẫn khóc, giọng vừa nghẹn ngào vừa đầy hận thù: ``Cô không hiểu đâu! Đó là bảo vật gia truyền! Nó có linh hồn! Và giờ thì\dots\ giờ thì\dots''

Cô không nói hết câu, chỉ có thể ôm đầu, lắc lắc trong sự tuyệt vọng. Choco nhìn cô, rồi nhìn bộ đàm, rồi lại nhìn cô, không biết nên ưu tiên việc gì trước: dỗ dành Ayame hay kiểm tra tình hình gỡ bom.

\ct

\pov{Hikawa Kuon}

Sợi dây màu đỏ đã bị cắt.

Kanata mắt trợn trừng, mặt cắt không còn giọt máu sau cú cắt đầy quyết đoán đó của Coco. Cô hét lên, giọng chói tai như thể vừa chứng kiến một tai nạn xe hơi: ``\textbf{Trời ạ! Bà cắt sợi đó làm gì vậy!? Làm cái gì!? Kuon-san đã bảo cắt dây xanh cơ mà!}''

Nàng Rồng vẫn giữ vẻ mặt hồn nhiên như thể vừa bấm nút ``thích'' trên Twitter mà không suy nghĩ gì. Cô nhìn sợi dây đỏ đã đứt, rồi quay sang chúng tôi, nở một nụ cười đầy tự hào: ``Tôi đã tìm hiểu kỹ trước khi làm đấy. Hôm qua tôi lên Yahoo Hỏi Đáp, tìm kiếm chơi `cách gỡ bom' và có một chuyên gia trả lời rằng: màu đỏ là dây của hệ thống an ninh. Nó chỉ có tác dụng báo động nếu cắt nhầm, nhưng nếu cắt đúng lúc thì nó sẽ tự ngắt. Còn màu xanh là dây nối trực tiếp với kíp nổ - cắt vào đó là chúng ta thành món khai vị luôn.''

Cô nói như thể đang giải thích một công thức nấu ăn, tay còn vẽ vẽ lên không trung. Tôi đứng hình.

Kanata thì ngược lại, phản ứng mạnh hơn nhiều. Cô ôm đầu, hào quang trên đầu xoay như điên: ``Bà điên rồi! Bà tin vào Yahoo Hỏi Đáp ư!? Đó là trang web mà ai cũng có thể viết! Có khi người trả lời là một đứa trẻ lớp ba cơ đấy! Sao bà không hỏi anh quản lý cho chắc?''

Coco vẫn điềm nhiên như không có chuyện gì xảy ra, nhìn nàng Thiên thần và tôi với ánh mắt đầy bí ẩn, rồi nói một câu khiến cả hai chúng tôi phải sởn gai ốc: ``Kanatan\dots\ Cô đã nghe đến câu `cái chết là sự giải thoát' chưa?''

Tôi bắt đầu toát mồ hôi, đưa tay lau trán: ``Em nói như thể chúng ta sẽ đi gặp ông bà ấy\dots\ Anh không muốn giải thoát ở đây đâu.''

Còn Kanata có vẻ phản ứng mạnh hơn tôi nhiều. Cô gần như nhảy dựng lên, chỉ tay vào mặt Coco: ``Bà hoảng quá nên mới đi cắt bừa thì có đó! Chứ không đời nào có người tỉnh táo lại tin vào câu trả lời trên Yahoo Hỏi Đáp trong một tình huống sinh tử! Mà Kuon-senpai này, anh học cái bài đó đâu ra vậy!?''

Một lần nữa, tôi đáp lại, giọng đều đều như đang đọc báo cáo: ``Kinh nghiệm quân sự thôi. Một tháng đi nghĩa vụ ở trường bắn. Bọn anh được học một chút về mìn và bom - dây xanh thường là dây an toàn, dây đỏ là dây phát nổ. Nhưng cũng có trường hợp ngược lại.''

Cả ba chúng tôi nhìn chằm chằm vào quả bom, trên đó chỉ còn lại hai mươi bảy giây, và không có dấu hiệu dừng lại. Mỗi giây như một tiếng búa đập vào ngực tôi.

Kanata hoảng loạn, nhảy lên đầu Coco, hét toáng lên: ``Đồng hồ vẫn chưa dừng lại kìa! Ta phải làm gì đây!?''

Coco vẫn bình thản như không có chuyện gì, lắc đầu nhẹ: ``Bình tĩnh nào! Nếu cắt đúng sợi thì sẽ vô hiệu được quả bom, chỉ là không vô hiệu được cái đồng hồ thôi! Nó vẫn chạy, nhưng không có gì xảy ra cả!''

Chúng tôi tròn mắt nhìn nhau với lập luận của em ấy: ``Thật à?''

Coco gật đầu chắc nịch, mặt vẫn tỉnh rụi: ``Ừ, trên Yahoo Hỏi Đáp có một chuyên gia khác nói vậy đấy!''

Tôi ngán ngẩm với lời giải mà em ấy tìm được. Tưởng em ấy tìm trên trang nào uy tín hơn chứ - như một trang quân sự uy tín, hoặc ít nhất là một video hướng dẫn từ MIT. Nhưng không, lại là Yahoo Hỏi Đáp.

Kanata nhún vai, coi như nhiệm vụ của chúng tôi đã xong: ``Sao cũng được! Vậy là ta gỡ được quả bom rồi nhỉ?''

Tôi nhìn chằm chằm vào quả bom vẫn đang đếm ngược kia, nhìn sợi dây xanh còn lại vẫn chưa cắt, và một cảm giác bất an vụt qua trong tôi: ``\dots Anh nghĩ là em nên cắt nốt cái dây kia đi cho an toàn. Linh cảm mách bảo vậy.''

Lập tức, Kanata liền thốt lên với vẻ hoảng sợ, như thể hét thẳng vào mặt tôi: ``\textbf{Anh bị điên rồi sao!? Anh muốn chúng ta banh xác chắc!? Vừa nãy cô bạn của em mới giải thích rằng dây xanh là dây nổ, giờ anh lại muốn cắt nốt!?}''

Coco bật cười, rồi nảy ra một sáng kiến bất ngờ, như một tia sáng lóe lên trong bóng tối: ``Hay là ta cứ cho nó vào tủ lạnh cho chắc đi!''

Tôi ngơ ngác với ý tưởng có phần ngớ ngẩn này, không biết nên cười hay nên khóc: ``\dots Hả? Tại sao chứ?''

``Có lẽ là\dots\ tuyết bên trong tủ đông sẽ dập quả bom ấy chăng?''  Coco đáp lại, lè lưỡi, trông như một đứa trẻ vừa đề xuất một trò chơi mới.

\dots Thật sự tôi cũng không biết phải giải thích thế nào với logic vốn dĩ đã như con số không của em ấy nữa. Tuyết dập bom? Tuyết làm giảm nhiệt độ, nhưng bom nổ hay không không liên quan đến nhiệt độ, mà liên quan đến dây điện và cơ chế kích nổ mà!

Thôi thì, nếu không nói được thì tốt nhất là thế này\dots

\pov{-}

Chưa đợi ai kịp phản ứng, nàng Thiên sứ - vẫn ngồi trên lưng Coco - vung tay ném thẳng quả bom vào tủ đông, rồi đóng sập cửa lại như thể đó là\dots\ một gói mì tôm ăn dở. Không một chút do dự, không một tiếng thở dài, chỉ có một tiếng ``cạch'' khô khốc của cánh cửa tủ đông khép lại.

\img{5-6i}

Khi quay đầu lại định báo cáo với Kuon, thì cậu đã biến mất từ lúc nào. Cô nhìn quanh, rồi hỏi với vẻ ngạc nhiên: ``À-rế? Kuon-paisen đâu rồi? Mới vừa nãy còn đứng đây cơ mà?''

Kanata buông một câu như đang dè bỉu cậu Tổng, giọng đầy tự mãn: ``Chắc là nhát quá nên chạy mất rồi. Đúng là gà! Mới có tí bom thôi mà đã chạy.''

Coco nhún vai, dường như đồng ý với suy nghĩ của nàng Thiên sứ. Cô khoanh tay, nhìn theo hướng Kuon đã biến mất, rồi thở dài: ``Anh ấy có vẻ không thích mạo hiểm. Nhưng cũng dễ hiểu thôi, quản lý mà, phải giữ mạng để còn ký giấy tờ nữa chứ.''

Ngay trước mặt hai người, Choco đang dắt theo Ayame tiu nghỉu như vừa bị mất sổ gạo. Nàng Oni cúi đầu, hai tay xấu hổ đan vào nhau, giọng lí nhí: ``Xin lỗi mấy đứa\dots\ achint\~{} Chị chỉ định đùa một chút thôi mà\dots''

Có vẻ như cô đã bị Choco giảng giải một bài dài về sự nguy hiểm của việc để bom lung tung, và đồng thời cũng nhận ra mối nguy hiểm khi để bom trong nhà. Choco đã chỉ cho cô thấy rằng những quả bom nhỏ có thể gây họa lớn hơn nhiều so với mấy con gián.

May mắn thay, nàng Thiên sứ lại tỏ ra vui vẻ, giơ ngón cái lên tỏ ý bỏ qua cho tiền bối, miệng cười tươi: ``Câu xoan-nỗi được chấp nhận! Dù sao cũng không có ai bị thương, quan trọng là mọi người vẫn còn sống!''

Và thế là, bốn con người ấy vui vẻ cười đùa với nhau, mọi thứ trở lại bình thường\dots

%==========================================TRUYỆN PHỤ=========================================
\omk

\pov{Hikawa Kuon}

\dots hoặc là không. Chắc chắn rồi. Vì khi bạn có một quả bom trong tủ đông, chẳng có gì là bình thường cả.

Vừa chạy thục mạng khỏi văn phòng - vì tôi đã ngửi thấy mùi nguy hiểm từ khi Coco đề xuất cho bom vào tủ lạnh - tôi lôi điện thoại ra, bấm số gọi cho một đứa bạn học cùng cấp ba của tôi: Morizaki Koheki (\ruby{森}{もり}\ruby{崎}{ざき}\ruby{個}{こ}\ruby{壁}{へき}, \ruby{Mô-ri}{Sâm}-\ruby{da-ki}{Kỳ}\ \ruby{Cô}{Cá}-\ruby{hê-ki}{Bích}).

Thằng này học giỏi Hóa, hiện tại đang đầu quân cho một học viện kỹ thuật quân sự ở nơi tôi ở. Chúng tôi vẫn thỉnh thoảng nói chuyện với nhau khi cậu ta có thời gian rảnh. Nó là người duy nhất tôi biết có thể giải thích mấy thứ như bom mìn mà không cần tra cứu mạng.

Điện thoại đổ chuông vài hồi, rồi một giọng nói lười biếng cất lên, như thể vừa tỉnh dậy sau một giấc ngủ trưa dài: ``\vie{A-lô? Ai đấy? Nếu là bán bảo hiểm thì tôi cúp máy đấy.}''

Tôi thở phào nhẹ nhõm khi đầu dây bên kia kết nối. Rất nhiều lần cậu bạn không bắt máy, hoặc là do đang ở trường quân sự, hoặc là không thèm trả lời. Tôi vừa nói vừa thở dốc: ``\vie{Là tao, Kuon đây. Tao hỏi mày một chút này\dots}''

Tôi bắt đầu kể lại toàn bộ hành trình gỡ bom của mình với Kanata và Coco. Tất nhiên, tôi không nhắc đến tên họ, vì kiểu gì thằng này cũng sẽ bảo tôi là một tên mê anime harem. Tôi chỉ nói là ``mấy đứa đồng nghiệp mới'' để đỡ mất mặt.

Sau khi thuật lại chi tiết loại bom mà chúng tôi vừa gỡ - từ hình dáng, màu sắc, đến hai sợi dây xanh đỏ - tôi nghe thấy một tiếng thở dài não nề từ đầu dây bên kia: ``\vie{Bọn kia cắt dây đỏ là dở cmn rồi. Xong rồi lại bỏ nó vô tủ lạnh? Thế càng dở hơn.}''

``\vie{Sao lại thế?}''  tôi nghiêng đầu, mặc dù biết rằng anh ta không thể thấy.

``\vie{Đây này-}''

{\Large *BÙM!*}

Chưa kịp để cậu ta giải thích, một tiếng nổ vang trời phát ra từ văn phòng, đủ mạnh để làm rung cả khu phố. Tôi bật ngửa ra sau vì choáng, suýt làm rơi điện thoại. Một luồng gió nóng thổi qua, mang theo mùi khét và bụi.

\img{5-6k}

Khói bốc lên nghi ngút từ tòa nhà, cùng với những gì còn sót lại từ văn phòng sau vụ nổ trời giáng đó. Tôi có thể thấy những mảnh kính vỡ bay ra từ cửa sổ, và một vài món đồ nội thất văng ra ngoài đường như thể đang chạy trốn.

``\vie{Cái gì vậy!?}''  Koheki ngạc nhiên, giọng đầy hoảng hốt sau vụ nổ đó. ``\vie{Có vụ nổ ở chỗ mày à? Mày có đang ở gần đó không?}''

Tôi ngập ngừng một lúc lâu, rồi tiếp tục trong tiếng thở dài, như thể vừa thú nhận một tội ác: ``\vie{\dots Quả bom tao vừa nói đấy. Nó vừa phát nổ.}''

``\dots''  một khoảng lặng dài ở đầu dây bên kia, đủ để tôi nghe thấy tiếng thở của Koheki. Rồi cậu ta nói, giọng đầy bất lực: ``\vie{Ủa, thế mày tưởng tao nãy giờ đang nói đùa sao?}''

``\vie{Mày nghĩ nãy giờ tao nói là đùa chắc?}''  tôi đáp lại với một cái thở dài. ``\vie{Nếu mà đùa thì tao đã chẳng phải làm phiền mày trong lúc đang nghỉ trưa giờ này.}''

Tôi hắng giọng một tiếng, rồi tiếp tục với chủ đề quả bom - giờ đã quá muộn để phá nó: ``\vie{Mà thế mày nói cắt dây đỏ không được là sao? Tại sao dây đỏ là dây giả?}''

Koheki thở dài, giọng như một giáo sư đang phải giải thích một điều hiển nhiên cho một học sinh mất gốc: ``\vie{Mày xem phim nhiều quá đấy. Trong thực tế, không có bất kỳ quy chuẩn màu sắc nào cho dây điện trong bom cả. Màu sắc chỉ là màu sắc, nó phụ thuộc vào những gì thằng chế tạo có sẵn.}''

``\vie{Ý mày là màu sắc không có nghĩa gì cả?}''

``\vie{Ờ. Có thể dây đỏ là dây nổ, dây xanh là dây giả, hoặc ngược lại. Hoặc cả hai đều là dây giả, và kíp nổ lại nằm ở chỗ khác. Hoặc tệ hơn, cả hai đều là dây nổ, và cắt bất kỳ sợi nào cũng khiến bom nổ. Không có sơ đồ mạch điện thì chẳng bố con thằng nào biết được bom hoạt động ra sao đâu.}''

Tôi ngẩn người, nhìn đống đổ nát của văn phòng: ``\vie{Vậy là\dots\dot{}việc cắt dây đỏ hay xanh không quan trọng?}''

``\vie{Chính xác. Cái mà bọn mày đã làm - cắt đại một sợi rồi ném bom vào tủ lạnh - đó là một trong những cách ngu ngốc nhất để xử lý bom. Mày có biết rằng trong thực tế, việc `cắt dây' gần như không bao giờ được sử dụng không? Người ta dùng robot, súng bắn nước áp lực cao, hoặc di dời bom đến nơi an toàn để nó phát nổ ở đó luôn.}''

``\vie{\dots Vậy mà phim Hollywood cứ làm quá lên.}''

``\vie{Ừ, phim Hollywood làm đếch gì có thật. Nó chỉ tạo kịch tính để bán vé thôi. Còn ngoài đời, nếu mày thấy một quả bom, mày gọi chuyên gia, đừng tự ý động vào. Và đừng bao giờ để mấy đứa đồng nghiệp mới của mày tự ý cắt dây dựa trên Yahoo Hỏi Đáp nữa.}''

Tôi gãi đầu một lúc lâu, nhìn bốn cô gái đang nằm lăn lóc trên đường, rồi cười trừ, vừa xấu hổ vừa hối hận: ``\vie{Lúc đầu tao bảo cắt dây xanh, nhưng Coco-- à nhầm, một đứa đồng nghiệp mới, ó bảo trên Yahoo Hỏi Đáp có người nói cắt dây đỏ mới đúng, rồi tự ý cắt luôn. Khi tao định cãi lại thì bị chửi là điên, vì dây xanh là dây nổ\dots\ Giờ thì tao biết là cả hai đứa đều điên.}''

Tôi có thể nghe thấy tiếng cậu ta đập tay vào trán vì sự ngớ ngẩn của tôi, hoặc là của Coco: ``\vie{Đúng là bọn mày óc chó thật\dots\ Lần sau gặp bom, mày gọi cho tao, đừng có tự làm chuyên gia. Còn gì để hỏi nữa không?}''

``\vie{Hết rồi.}''

``\vie{Ờ, thế nhá. Tao cũng chuẩn bị đi tập quân sự rồi đây.}''

``\vie{Ờ, tập vui vẻ. Khi nào rảnh em mình làm một bữa.}''

Tôi cúp máy, ngước mặt lên bầu trời xanh hòa với cột khói đen nghi ngút. Những mảnh giấy vụn bay trong gió, trông như tuyết trong một ngày đông - nhưng mà tuyết màu xám.

Bất chợt, tôi nhìn thấy bốn bóng người - Coco, Kanata, Ayame và Choco - những người đang rơi tự do trên cao xuống, tay chân vung vẫy trong không trung như bốn con búp bê vải bị ai đó ném đi. Có vẻ như vụ nổ đã đẩy họ ra khỏi tòa nhà.

Họ tiếp đất bằng một pha va chạm xuống mặt đất, toàn thân đen thui như bị nướng chín, quần áo bốc khói, mắt mất tiêu đồng tử. Tất cả họ đồng loạt rên la trong đau đớn, như một dàn hợp xướng thất bại.

Tôi chỉ biết khoanh tay thở dài, lắc đầu bất lực, nhìn họ nằm sóng soài trên vỉa hè: ``\textit{Bảo rồi không nghe cơ\dots\ Giờ thì họ sẽ phải đền một núi tiền đấy\dots\ Cửa sổ, bàn ghế, tủ đông, và cả mấy cái búp bê mà Ayame cứ khóc hoài\dots}''

\il

Từ trên trời rơi xuống thêm một vật thể lạ. Một mảnh giấy.

Nó nhẹ nhàng đáp xuống mặt đất ngay trước mặt tôi, trông có vẻ là một lá thư. Khói bụi vẫn còn bao phủ, nhưng mảnh giấy trắng tinh như thể không hề bị ảnh hưởng bởi vụ nổ. Thật kỳ lạ.

\begin{wrapfigure}[11]{l}{6cm}
    \includegraphics[width=8cm]{../../../../Image/Book5/5-6l.png}
\end{wrapfigure}

Tôi nhặt nó lên, phủi bụi. Không có ghi người gửi hay người nhận. Chỉ có một hình con cừu nhỏ in trên đó - một chú cừu đen nền vàng, trông như ai đó đã vẽ bằng tay.

Và ngay bên dưới là ba chữ to đùng, viết bằng mực đỏ: ``\textbf{Thư thách đấu.}''

Ngay khoảnh khắc đó, tôi nhắm mắt thở dài.

Không cần mở ra, tôi cũng biết nó là của ai gửi, và gửi cho ai.

Thư của ``con cừu''.

Và ``con cừu'' đó không ai khác ngoài\dots\ Tôi sẽ không nói tên của người đó ra - vì tôi biết kiểu gì cô ta cũng sẽ đọc được chỗ này và cười khoái chí - nhưng tôi chắc các bạn cũng đoán được đó là ai rồi. Cô ấy là người duy nhất trong \hll\ có biệt danh ``con cừu'' và thích gửi thư thách đấu kiểu này.

Tôi có thể đoán trước được nội dung bên trong mà không cần đọc, có thể là\dots

Một lời tuyên chiến đầy khí thế, với giọng văn kiểu: ``Hỡi các chiến binh, hãy chuẩn bị đi! Vì ta sẽ đến thử thách các ngươi vào tuần sau!''

Hoặc là, một vài câu chế giễu nhẹ nhàng (có thể là nhắm vào tôi) kiểu: ``Còn nhớ lần trước cậu thua tôi vụ đó không, Kuon-senpai? Haha, lần này tôi sẽ không để cậu có cơ hội gỡ gạc đâu!''

Mặc dù tôi còn chẳng nhớ tôi cá cược gì với họ. Có lẽ tôi đã quên vì nó không quan trọng, hoặc có lẽ tôi đã cố tình quên vì thua cuộc.

Hoặc cũng có thể là một dấu ấn riêng của người đó - một tựa đề đầy hoành tráng kiểu: ``Đệ Nhất Kiếm Cừu Đến Rồi!''

\dots Có lẽ vậy. Dù gì thì, đồng nghĩa với việc\dots

Lại có thêm một thành viên nữa của \textit{holoForce} sẽ đến văn phòng chúng tôi vào tuần sau, ngay sau khi nơi này được xây dựng lại. Tôi nhìn quanh đống đổ nát, những mảnh kính vỡ, và bốn cô gái đang nằm lăn lóc trên đường, rồi thở dài.

Hầy\dots\ chưa gì mà đã thấy mệt rồi đấy\dots

\begin{center}
    {\Large \abv{Bonus}}
\end{center}

\begin{figure}[H]
    \centering
    \includegraphics[width=12.5cm]{../../../../Image/Book5/5-6m.png}
\end{figure}

\end{document}